\paragraph{Letra A}

Usamos novamente o modelo paralelo para estimar os parâmetros da planta. Neste caso escolhemos um sistema simples com uma Resistência R em série com um Indutor L. A equação que descreve o sistema neste caso é $V = R*I + L\dot{I}$ que manipulamos para a forma $\dot{x} = -a*x + b*u$, onde $x=I; U=V;a = R/L;b = 1/L$. Os valores escolhidos foram de $R=12\Omega; L=1H$ para a letra A. Podemos notar que apesar de a saída do modelo da planta passar a ter o mesmo valor que a saída da planta após um certo intervalo de tempo, os valores dos parâmetros estimados não necessariamente irão convergir para os valores corretos.

\includegraphics{./pic/Q2letraAmontagem.png}
 
Podemos observar o comportamento da estimativa de R/L no gráfico a seguir:

\includegraphics{./pic/Q2letraAteta1.png}

E da estimativa de 1/L no gráfico a seguir:

\includegraphics{./pic/Q2LetraAteta2.png} 

E a saída do modelo seguindo a saída da planta (erro de saída convergindo para 0):

\includegraphics{./pic/Q2letraAe0.png}
 
\paragraph{Letra B}

Para a letra B foram escolhidos valores de $R=1.2\Omega e L=0.01H$, neste caso, como um sinal senoidal é persistentemente excitante para uma planta de ordem 1 podemos observar a convergência dos valores das estimativas de R/L e 1/L para os valores corretos.

\includegraphics{./pic/Q2letraBmontagem.png}
 
Podemos observar o comportamento da estimativa de R/L no gráfico a seguir:

\includegraphics{./pic/Q2letraBteta1.png} 

E o comportamento da estimativa de 1/L no gráfico a seguir:

\includegraphics{./pic/Q2letraBteta2.png} 

E a saída do modelo seguindo a saída da planta (erro de saída convergindo para 0):

\includegraphics{./pic/Q2letraBe0.png}
 
